\documentclass[10pt,a4paper]{amsart}
\usepackage[no-math]{fontspec}
\usepackage{amsmath,amssymb,mathtools,xfrac,xcolor}
\usepackage[scr=rsfs,cal=euler]{mathalfa}
\usepackage[a4paper,margin=22mm]{geometry}
\setmainfont{Latin Modern Roman}[SmallCapsFont=lmromancaps10-regular.otf]
\usepackage[colorlinks=true,linkcolor=blue,citecolor=teal,urlcolor=blue]{hyperref}
\sloppy\newcommand{\cf}{cf.\ }\newcommand{\resp}{resp.\ }
\multlinegap0pt
\let\epsilon\varepsilon

\hyphenation{authors rela-tion restric-tion}

\let\ts\textstyle
\newcommand\mto{\mathchoice{\longmapsto}{\mapsto}{\mapsto}{\mapsto}}

\newcommand{\Psfrac}[2]{(\sfrac{#1}{#2})}
\newcommand{\psfrac}[2]{\sfrac{(#1)}{#2}}
\newcommand{\spfrac}[2]{\sfrac{#1}{(#2)}}
\newcommand{\pspfrac}[2]{\sfrac{(#1)}{(#2)}}

\newcommand{\app}{\mathrm{app}}
\newcommand{\loc}{\mathrm{loc}}
\newcommand{\st}{\mathrm{st}}

\def\Xint#1{\mathchoice
{\XXint\displaystyle\textstyle{#1}}%
{\XXint\textstyle\scriptstyle{#1}}%
{\XXint\scriptstyle\scriptscriptstyle{#1}}%
{\XXint\scriptscriptstyle\scriptscriptstyle{#1}}%
\!\int}
\def\XXint#1#2#3{{\setbox0=\hbox{$#1{#2#3}{\int}$ }
\vcenter{\hbox{$#2#3$ }}\kern-.6\wd0}}
\def\ddashint{\Xint=}
\def\dashint{\Xint-}

\renewcommand{\div}{\operatorname{div}}

\theoremstyle{plain}
\newtheorem{thm}{Theorem}[section]
\newtheorem{prop}[thm]{Proposition}
\newtheorem{coro}[thm]{Corollary}
\newtheorem{lemme}[thm]{Lemma}

\theoremstyle{definition}
\newtheorem{defi}[thm]{Definition}
\newtheorem{rem}{Remark}[section]

\let\pa\partial
\let\na\nabla
\let\eps\varepsilon

\newcommand{\eff}{\mathrm{eff}}
\newcommand{\sym}{\operatorname{sym}}
\newcommand{\ssym}{\operatorname{ssym}}
\newcommand{\Lip}{\mathrm{Lip}}
\newcommand{\supp}{\operatorname{supp}}
\newcommand{\mO}{\mathcal{O}}
\newcommand{\mM}{\mathcal{M}}
\newcommand{\mS}{\mathcal{S}}
\newcommand{\mV}{\mathcal{V}}
\newcommand{\R}{\mathbb{R}}
\newcommand{\T}{\mathbb{T}}
\newcommand{\N}{\mathbb{N}}
\newcommand{\Z}{\mathbb{Z}}

\begin{document}
\title[Half-space Stokes depletion-layer expansion]{Arbitrary-order depletion-layer expansion for the continuous half-space Stokes model}
\author{David Gérard-Varet\and Amina Mecherbet}
\date{}\maketitle\section{Statement and required original model}
for some unknown translational and angular velocities $v_i, \omega_i \in \R^3$. The two relations in the last line, correspond to balances of forces and torques, and involve the Newtonian stress tensor $\Sigma_\mu(u,p) = 2\mu D(u) - p I_3$ with $D(u)=\frac{1}{2}(\nabla u + \nabla^\top u)$ the symmetric gradient of $u$. The constants $v_i, \omega_i$ can be seen as Lagrange multipliers associated with these constraints. The term $-mge$, $e = (0,0,1)^t$ encodes gravity and buoyancy effects: denoting $\rho_p$ and $\rho_f$ the particle and fluid masses per unit volume, one has $m = \frac{4}{3} \pi R^3 (\rho_p - \rho_f)$, while $g$ is the usual gravitational constant. This is the term responsible for sedimentation.
At first sight, these two phenomena may seem contradictory (downward {\it vs} upward flow just outside the depletion layers). One ambition of this paper is to clarify these different behaviors. More globally, we~wish here to analyze rigorously the effect of a depletion layer on a suspension, starting from the simplest model \href{https://jep.centre-mersenne.org/item/10.5802/jep.296.pdf#page=3}{(1)}. For $\eps > 0$ we consider the half-space $\Omega = \Omega^\eps = (-\eps,+\infty) \times \R^2$, where the depletion layer $D_\eps = (-\eps,0) \times \R^2$, of width $\eps$, is free of particles centers. Namely,
\begin{equation} \label{eq_assumption}
X_1, \dots, X_N \in K \Subset \overline{\Omega^0}, \quad \Omega^0 = \R^*_+ \times \R^2.
\end{equation}
where $\Sigma(u,p) = 2D(u) - p I$. Note that we could have considered some variation of~\href{https://jep.centre-mersenne.org/item/10.5802/jep.296.pdf#page=5}{(4)} with additional forcing $\int_{B_i} f$ and torque $\int_{B_i} f \times (x-X_i)$ included in the last two relations.
\setcounter{equation}{4}
\begin{equation} \label{limit_equation}
\begin{aligned}
-\Delta u^\eps + \na p^\eps & = f - \rho e \quad \text{ in } \: \Omega^\eps, \\
\div u^\eps & = 0 \quad \text{ in } \: \Omega^\eps, \\
u^\eps \vert_{\pa \Omega^\eps} & = 0.
\end{aligned}
\end{equation}
We assume in all what follows that $f \in L^1(\Omega^\eps) \cap L^\infty(\Omega^\eps)$. In~particular, $f \in L^{6/5}(\Omega^\eps)$ so that \href{https://jep.centre-mersenne.org/item/10.5802/jep.296.pdf#page=5}{(4)}, \resp \eqref{limit_equation}, has a unique solution $u^{N,\eps}$, \resp $u^\eps$ in
\[
V = \{ u \in H^1_{\loc}(\overline{\Omega^\eps}), \na u \in L^2(\Omega^\eps)\}.
\]
\setcounter{thm}{1}
\begin{equation} \label{crude_approx}
\begin{aligned}
-\Delta u^0 + \na p^0 &= f - \rho e \quad \text{ in } \: \Omega, \\
\div u^0 & = 0 \quad \text{ in } \: \Omega^0, \\
u^0 \vert_{\pa \Omega^0} & = 0,
\end{aligned}
\end{equation}
where the impact of the layer is totally neglected. In~this regard, we~will prove in Section \ref{sec_wall_law} the following theorem:
\begin{thm} \label{thm2}
Assume that $f$ and $\rho$ are smooth, compactly supported in $\overline{\Omega^0}$. Then, for all $m \in \N$, one has
\[
u^\eps = u^0 + \eps u^1 + \dots + \eps^m u^m + O(\eps^{m+1}) \quad \text{ in } \: \dot{H}^1(\Omega^0),
\]
where $u^0$ solves \eqref{crude_approx}, and for each $i \ge 1$, $u^i$ solves an homogeneous Stokes equation in~$\Omega^0$ with inhomogeneous Dirichlet data at $\pa \Omega^0$ coming from lower order profiles.
\end{thm}
This theorem provides an asymptotic expansion of $u^\eps$ at arbitrary order, where each successive correction refines the description of the effect of the depletion layer. One could relax the assumptions on the smoothness of $\rho$ and $f$: for instance, it is enough that $\na u^0,p^0$ belong to $H^{m+2}(\Omega^0)$ to have an expansion with $m$ terms.
\setcounter{section}{3}
\section{The approximate continuous model} \label{sec_wall_law}The purpose of this section is to analyze system \eqref{limit_equation}, and to prove Theorem \ref{thm2} and Corollary \href{https://jep.centre-mersenne.org/item/10.5802/jep.296.pdf#page=7}{1.3}. We assume here that $f$ and $\rho$ are smooth and compactly supported in~$\overline{\Omega^0}$.

\subsection{Asymptotic expansion}
We will construct an approximation of the solution $(u^\eps,p^\eps)$ of \eqref{limit_equation}, of the following form:
\begin{equation*}
\left\{
\begin{aligned}
(u^\eps_{\app},p^\eps_{\app})(x_1,x_2,x_3) & \approx \sum_{i=0}^m \eps^i (u^i, p^i)(x_1,x_2,x_3), \quad x = (x_1,x_2,x_3) \in \Omega^0, \\
(u^\eps_{\app}, p^\eps_{\app}(x_1,x_2,x_3) & \approx \sum_{i=0}^m \eps^i (U^i, P^i)(x_1/\eps,x_2,x_3), \quad x = (x_1,x_2,x_3) \in \Omega^\eps\setminus\Omega^0.
\end{aligned}
\right.
\end{equation*}
Note that this approximation is made of two parts: an interior one, in $\Omega^0$, with a regular expansion in $\eps$, and a boundary layer part, localized in the depletion layer $\Omega^\eps \setminus\nobreak \Omega^0$. The boundary layer expansion involves boundary layer profiles $U^i = U^i(s,x_2,x_3)$, with the variable $s \in (-1,0)$ that stands for
$x_1/\eps$. It is also convenient to set:
\[
(u^i, p^i)= 0, \quad (U^i, P^i) = 0 \quad \text{ for } \: i < 0.
\]
Plugging the interior expansion in the Stokes equation, we~find: for all $i \ge 0$, in $\Omega^0$
\begin{equation} \label{stokesint}
\begin{aligned}
-\Delta u^i + \na p^i & = \delta_{0i} \big(f - \rho e \big), \\
\div u^i & = 0.
\end{aligned}
\end{equation}
Plugging the boundary layer expansion in the Stokes equation, we~find: for all $i \ge 0$, for all $(s,x_2,x_3) \in (-1,0) \times \R^2$:
\begin{align}
-\pa^2_s U^i_1 - (\pa^2_2 + \pa^2_3) U^{i-2}_1 + \pa_s P^{i-1} & = 0, \label{stokesx} \\
-\pa^2_s U^i_2 - (\pa^2_2 + \pa^2_3) U^{i-2}_2 + \pa_2 P^{i-2} & = 0, \label{stokesy}\\
-\pa^2_s U^i_3 - (\pa^2_2 + \pa^2_3) U^{i-2}_3 + \pa_3 P^{i-2} & = 0, \label{stokesz} \\
\pa_s U^i_1 + \pa_2 U^{i-1}_2 + \pa_3 U^{i-1}_3 & = 0 \label{stokesdiv}
\end{align}
The Dirichlet condition $u^\eps_{\app}\vert_{\pa \Omega^\eps} \approx 0$ further yields: for all $i \ge 0$,
\begin{equation} \label{dirichlet}
U^i\vert_{s = -1} = 0.
\end{equation}
Eventually, conditions
\[
[u^\eps_{\app}]\vert_{\pa \Omega^0} \approx 0, \quad [\Sigma(u^\eps_{\app},p^\eps_{\app})n]\vert_{\pa \Omega^0} \approx 0,
\]
which reflect continuity of the velocity field and the stress tensor at the interface $\pa \Omega^0$ give for all $i \ge 0$:
\begin{align}
U^i\vert_{s=0} & = u^i\vert_{x_1=0} \label{continuity_vel} \\
\pa_s U^i \vert_{s=0} - P^{i-1}\vert_{s=0} (1,0,0)^t & = \pa_1 u^{i-1}\vert_{x_1=0} - p^{i-1}\vert_{x_1=0} (1,0,0)^t. \label{continuity_stress}
\end{align}

\subsubsection*{Computation of the first terms.} We compute $U^0, u^0,p^0$. First, from \eqref{stokesdiv} and \eqref{dirichlet}, we~get $\pa_s U^0_1 = 0$ and $U^0_1\vert_{s=-1} = 0$, which implies $U^0_1 = 0$. Then, from \eqref{stokesy}, \eqref{dirichlet} and~\eqref{continuity_stress}, we~find
\[
\pa^2_s U^0_2 = 0, \quad U^0_2\vert_{s=-1} = 0, \quad \pa_s U^0_2\vert_{s=0} = 0,
\]
which leads to $U^0_2 = 0$. Similarly, $U^0_3 = 0$. Hence, $U^0 = 0$. Next, we~consider \eqref{stokesint} and~\eqref{continuity_vel} at rank $i=0$. They imply that $(u^0,p^0)$ satisfies \eqref{limit_equation}. In~particular, as~$f$ belongs to $L^{6/5}(\Omega_0) \cap H^m(\Omega^0)$, $(\na u^0,p^0)$ belongs to $H^m(\Omega^0)$ for all $m$.

\subsubsection*{Computation of next order terms.} We now assume that $i \ge 1$ and that profiles $U^k$, $P^{k-1}$, $u^k$, $p^k$ are known for all $k \le i-1$, with:
\[
U^k, P^{k-1} \in H^\infty((-1,0) \times \R^2), \quad \na u^k, p^k \in H^\infty(\Omega^0).
\]
We now show how to construct $U^i,P^{i-1}, u^i, p^i$. Considering \eqref{stokesdiv} and \eqref{dirichlet} yields
\[
U^i_1(s,\cdot) = - \int_{-1}^s (\pa_2 U^{i-1}_2 + \pa_3 U^{i-1}_3)(t,\cdot) dt \in H^\infty((-1,0) \times \R^2).
\]
Next, we~consider \eqref{stokesx} and \eqref{continuity_stress} which allow the calculation of \hbox{$P^{i-1} \!\!\in\! H^\infty(\!({-}1,0)\!\times\! \R^2)$}:
\[
P^{i-1}(s,\cdot) = \pa_s U^i_1\vert_{s=0} - \pa_1 u^{i-1}_1\vert_{x_1=0} + p^{i-1}\vert_{x_1=0}
+ \int_0^s \big(\pa^2_s U^i_1 + (\pa^2_2 + \pa^2_3) U^{i-2}_1 \big)(t,\cdot) dt.
\]
Next, we~consider \eqref{stokesy}--\eqref{dirichlet}--\eqref{continuity_stress}: these relations imply
\[
U^i_2(s,\cdot) = - \int_{-1}^s \int_{0}^t \Big((\pa^2_2 + \pa^2_3) U^{i-2}_2 + \pa_2 P^{i-2} \Big)(t',\cdot) dt' dt + (s+1) \pa_1 u^{i-1}_2\vert_{x_1=0}.
\]
A similar expression holds for $U^i_3$. Both $U^i_2, U^i_3 \in H^\infty((-1,0) \times \R^2)$. Eventually, we~see by \eqref{stokesint} and \eqref{continuity_vel} that $(u^i,p^i)$ solve a homogeneous Stokes equation with the inhomogeneous Dirichlet condition
\[
u^i\vert_{x_1=0} = U^i\vert_{x_1=0}.
\]
By standard regularity results for the Stokes equation, \cf
\cite[Th.\,IV.3.2 \& IV.3.3]{Galdi}, $\na u^i, p^i \in H^\infty(\Omega^0)$.

\subsection{Proof of Theorem \ref{thm2}}
We will show here that the expansion constructed in the previous section provides an approximate solution of the true solution $u^\eps$ of \eqref{limit_equation}. Given $m \in \N$, we~first introduce $\tilde{u}^{m+2}$ the solution in $\dot{H}^1(\Omega^\eps)$ of the Stokes problem
\begin{equation*}
\left\{
\begin{aligned}
-\Delta \tilde{u}^{m+2} + \na \tilde{p}^{m+2} & = 0 \quad \text{in } \Omega^\eps, \\
\div \tilde{u}^{m+2} & = - \pa_1 u^{m+2}_1 1_{\Omega^0} \quad \text{in } \Omega^\eps, \\
\tilde{u}^{m+2}\vert_{\pa \Omega^\eps} & = 0.
\end{aligned}
\right.
\end{equation*}
Such $\tilde{u}^{m+2}$ exists, \cf \cite[Th.\,IV.3.3]{Galdi}, and satisfies
\[
\|\na \tilde{u}^{m+2}\|_{L^2(\Omega^\eps)} \lesssim \|\pa_1 u^{m+2}_1 1_{\Omega^0} \|_{L^2(\Omega^\eps)} \lesssim 1.
\]
We then define $(u^\eps_{\app}, p^{\eps}_{\app})$ as follows:
\begin{equation*}
\begin{aligned}
(u^\eps_{\app},p^\eps_{\app})(x,y,z) & = \sum_{i=0}^{m+1} \eps^i (u^i, p^i)(x,y,z) + \eps^{m+2} (u^{m+2}_1,0,0,0)(x,y,z) \\
& \hspace*{1cm}+ \eps^{m+2} (\tilde{u}^{m+2},\tilde{p}^{m+2})(x,y,z), \qquad X = (x,y,z) \in \Omega^0, \\
(u^\eps_{\app}, p^\eps_{\app})(x,y,z) & = \sum_{i=0}^{m+1} \eps^i (U^i, P^i)(x/\eps,y,z) + \eps^{m+2} (U^{m+2}_1,0,0,0)(x/\eps,y,z) \\
& \hspace*{1cm}+ \eps^{m+2} (\tilde{u}^{m+2},\tilde{p}^{m+2})(x,y,z), \qquad \: X = (x,y,z) \in \Omega^\eps\setminus\Omega^0.
\end{aligned}
\end{equation*}
We then introduce $v^\eps = u^\eps - u^\eps_{\app}$, $q^\eps = p^\eps - p^\eps_{\app}$, solving
\begin{align}
-\Delta v^\eps + \na q^\eps & = - \eps^{m+2} (\Delta u^{m+2}_1,0,0) \quad \text{ in } \: \Omega^0,\notag \\
-\Delta v^\eps + \na q^\eps & = - \sum_{i=m}^{m+1} \eps^i (\pa^2_2 + \pa^2_3) (\eps U^{i+1}_1,U^i_2,U^i_3) - \sum_{i=m}^{m+1} \eps^i (0, \pa_2 P^i, \pa_3 P^i) \quad \text{ in } \: D^\eps, \notag \\
\div v^\eps & = 0 \quad \text{ in } \: \Omega^\eps, \\
[v^\eps]\vert_{\pa \Omega^0} & = 0, \notag\\
[\Sigma(v^\eps,q^\eps)n]\vert_{\pa \Omega^0} & = - \eps^{m+1}(\eps \pa_1 u^{m+2}_1\vert_{x_1=0},\pa_1 u^{m+1}_2\vert_{x_1=0}, \pa_1 u^{m+1}_3\vert_{x_1=0}), \notag \\
v^\eps \vert_{\pa \Omega^\eps} & = 0. \notag
\end{align}
Performing a simple energy estimate, we~find
\begin{align*}
\int_{\Omega^\eps} |\na v^\eps|^2 = \eps^{m+2} \int_{\Omega^0} \na u_1^{m+2} : \na v_1^\eps - \int_{D^\eps} R^\eps \cdot v^\eps + \int_{\pa \Omega^0} r^\eps \cdot v^\eps,
\end{align*}
where
\begin{align*}
R^\eps & = - \sum_{i=m}^{m+1} \eps^i (\pa^2_2 + \pa^2_3) (\eps U^{i+1}_1,U^i_2,U^i_3) - \sum_{i=m}^{m+1} \eps^i (0, \pa_2 P^i, \pa_3 P^i) = O(\eps^m), \\
r^\eps & = \eps^{m+1}(\eps \pa_1 u^{m+2}_1\vert_{x_1=0},\pa_1 u^{m+1}_2\vert_{x_1=0}, \pa_1 u^{m+1}_3\vert_{x_1=0}) = O(\eps^{m+1}).
\end{align*}
By the Poincaré and trace inequalities
\[
\| v^\eps \|_{L^2(D^\eps)} \lesssim \eps \|\na v^\eps \|_{L^2(D^\eps)}, \quad \| v^\eps \|_{L^2(\pa\Omega^0)} \lesssim \sqrt{\eps} \|\na v^\eps \|_{L^2(D^\eps)},
\]
we infer easily
\[
\|\na v^\eps \|_{L^2(\Omega^\eps)} \lesssim \eps^{m+1}.
\]
Restricting to $\Omega^0$, Theorem \ref{thm2} follows.


\begin{thebibliography}{10}
\bibitem[10]{Galdi} G. P. Galdi – An introduction to the mathematical theory of the Navier-Stokes equations. Steady-state problems, second ed., Springer Monographs in Math., Springer, New York, 2011. \url{https://doi.org/10.1007/978-0-387-09620-9}
\end{thebibliography}
\end{document}
